\documentclass{article}
\usepackage[T1]{fontenc}
\usepackage{lmodern}
\usepackage{graphicx}
\usepackage{amsmath,amssymb}
\usepackage[a4paper,margin=2cm]{geometry}
\usepackage[absolute,overlay]{textpos}
\setlength{\TPHorizModule}{1mm}
\setlength{\TPVertModule}{1mm}
\usepackage[indonesian]{babel}
\usepackage{hyperref}
\setlength{\emergencystretch}{2em}
% unit: Halaman 11

\begin{document}

% === Halaman 11 (python/native) ===

11

Operasi Bitwise XOR

• Jika dua rangkaian dioperasikan dengan XOR, maka 

operasinya dilakukan dengan meng-XOR-kan setiap bit yang 

berkoresponden dari kedua rangkaian bit tersebut. 


Contoh: 10011  11001 = 01010 

yang dalam hal ini, hasilnya diperoleh sebagai berikut: 




1 

   0      0       1         1 




1 

   1      0       0         1     


 1  1   0  1  0  0  1 0  1  1  



0 

  1 

     0       1         0

\end{document}
